\documentclass[12pt]{article}
\usepackage[utf8]{inputenc}
\usepackage[T1]{fontenc}
\usepackage[margin=1in]{geometry}
\usepackage{setspace}
\usepackage{xcolor}
\usepackage[colorlinks=true, urlcolor=blue!60!black, linkcolor=blue!60!black]{hyperref}
\setlength{\parindent}{0em}
\setlength{\parskip}{0.45em}
\renewcommand{\baselinestretch}{1.0}
\begin{document}
\begin{center}
{\Large \textit{Haaland v. Brackeen}, 599 U.S. 255 (2023)\par}
\vspace{0.3em}
{\normalsize Case Brief --- One-Page Edition\par}
\end{center}

\textbf{Memory Jogger.} ICWA upheld against Article~I and anticommandeering attacks; equal protection claims dismissed for lack of standing (\href{https://supreme.justia.com/cases/federal/us/599/21-376/}{599 U.S. 255}).

\textbf{Facts.} ICWA requires ``any party'' seeking to remove an Indian child to show ``active efforts'' to keep the family together and prefers extended family, tribal members, and ``other Indian families'' in adoptions. The Brackeens and other non-Indian families, joined by Texas, challenged it after disputes over the adoptions of A.L.M., Baby O., and Child P.

\textbf{Procedural History.} The district court struck down ICWA; the en banc Fifth Circuit upheld Congress's power but invalidated several provisions.

\textbf{Votes.} 7--2. Barrett wrote for the Court; Gorsuch and Kavanaugh concurred; Thomas and Alito dissented.

\textbf{Holding.} ICWA is within Congress's Indian-affairs power, with no ``family law carveout''; it does not commandeer the States; and neither the families nor Texas had standing to raise equal protection.

\textbf{Reasoning.} The power rests on the Indian Commerce Clause, treaties, structure, and the trust relationship. Section 1912 binds ``any party,'' \S~1915 imposes no search duty on state agencies, and recordkeeping is ancillary. State judges who apply the preferences are not bound by a judgment against federal officials, and Texas has no equal protection rights. Kavanaugh called the equal protection issue ``serious''; Thomas and Alito saw plenary power unmoored from the text.

\textbf{Examples.} \emph{Same-Side:} an agency suing the Secretary lacks standing. \emph{Opposite-Side:} a state-only search duty. \emph{Fence-Sitter:} equal protection raised in state court.

\textbf{Critique.} The Court never defined the borders of an ``unwieldy'' doctrine, and its ``any party'' test elevates form over substance where state agencies bring most cases; Gorsuch and Thomas reached opposite originalist answers.

\textbf{Key Quotation.} ``It is a federal court's judgment, not its opinion, that remedies an injury; thus it is the judgment, not the opinion, that demonstrates redressability. The individual petitioners can hope for nothing more than an opinion, so they cannot satisfy Article III.'' --- Justice Barrett, 599 U.S. 255.

\end{document}
